% Supplementary material for main.tex.  Build main.tex first, then this file, then main.tex again
% (labels are exchanged through the .aux files with the xr package).
\documentclass[conference,onecolumn,11pt]{IEEEtran}
\IEEEoverridecommandlockouts
\usepackage{cite}
\usepackage{amsmath,amssymb,amsfonts}
\usepackage{graphicx}
\usepackage{textcomp}
\usepackage{url}
\usepackage{xcolor}
\usepackage{booktabs}
\usepackage{multirow}
\usepackage{array}
\usepackage{enumitem}
\usepackage{xr}
\externaldocument[M-]{main}
\newcommand{\E}{\mathbb{E}}
\newcommand{\Prob}{\mathbb{P}}
\graphicspath{{figures/}}
\setlength{\emergencystretch}{3em}
\input{tables/macros}
\renewcommand{\thetable}{S\arabic{table}}
\renewcommand{\thefigure}{S\arabic{figure}}
\renewcommand{\thesection}{S\arabic{section}}
\renewcommand{\thesubsection}{\thesection.\arabic{subsection}}

\begin{document}

\title{Supplementary Material for ``Verified Stopping Dominates Adaptive Allocation in Multi-Container DNA
Storage: A Computational Study on a Digital Twin and Real Read Streams''}

\author{\IEEEauthorblockN{Sai Mouli}
\IEEEauthorblockA{saimouli12345678@gmail.com}
\and
\IEEEauthorblockN{J Geetha}
\IEEEauthorblockA{geetha@msrit.edu}
\and
\IEEEauthorblockN{Pruthviraj M. Savanur}
\IEEEauthorblockA{rajsavanur2003@gmail.com}
}

\maketitle

\noindent Section, table and figure numbers without the prefix ``S'' refer to the main paper.

%=====================================================================
\section{Positioning against prior work}\label{supp:lit}
%=====================================================================
Table~\ref{tab:literature_comparison} compares this work with the closest prior work discussed in
Section~\ref{M-sec:related}.
\input{tables/literature_comparison}

%=====================================================================
\section{Operational specification of the retrieval loop}\label{supp:alg}
%=====================================================================
Table~\ref{tab:algorithms} specifies the loop and every allocation rule: initialisation, batch size,
tie-breaking, when the posterior is updated, the stopping test and how an exhausted stream is handled.
The two model-free baselines of Table~\ref{M-tab:r9base}(b) are specified in Section~\ref{M-sec:res-extra}.
\input{tables/algorithms}

%=====================================================================
\section{Information-greedy allocation}\label{supp:info}
%=====================================================================
Information-greedy reads the container with the largest posterior variance of $\ln\theta_j$ until every
variance falls below $10^{-3}$, then switches to coverage-greedy.  Under the primary container model the
near-dead container's variance never falls below the threshold, so the rule spends 66\% of its reads
there, against 16\% for round-robin, and costs 441\% more than round-robin on the baseline.  A richer
model gives it more uncertainty to chase: with the dropout model the penalty grows to 612\%.  In the
homogeneous control it is indistinguishable from round-robin, and after an outright failure it is 81\%
costlier on the retrievals that succeed.  Table~\ref{tab:suppinfo} collects every result for this rule.
It is a pure information-gathering rule with no cost term, and we report it for completeness rather than
as a competitive baseline.
\input{tables/supp_infogreedy}

%=====================================================================
\section{One failed container: conditional and unconditional cost}\label{supp:uncond}
%=====================================================================
\input{tables/r4_uncond}
With one failed container 34\% of retrievals succeeded, and the success probability was identical for
every policy and every static plan (Table~\ref{tab:uncond}): whether the file can be recovered is decided
by which fragments survive in the five remaining containers, not by the reading order.  Among the 102
retrievals that succeeded, UCB used 15.49k reads against round-robin's 20.31k, 24\% fewer.  Counting all
300 attempts, including the 66\% that read the whole 72k-read stream and failed, UCB used 52.79k against
round-robin's 54.42k, a 3\% reduction.  The two numbers answer different questions: a user who retries
until success pays the conditional number, and a user who pays per attempt pays the unconditional one.
After an outright failure, posterior rules therefore cut the cost of the retrievals that succeed by up to
24\% without improving the probability of success, and change the cost per attempt by 3\%.

%=====================================================================
\section{Calibration of the posterior and re-selection of $\rho_U$}\label{supp:calib}
%=====================================================================
\input{tables/r4_calib}
On the ten selection replicates, the 90\% credible-interval coverage of $\theta^\ast$ rose with $\rho_U$
from 0.45 (binomial) through 0.63, 0.73 and 0.80 to 0.87 at $\rho_U{=}0.1$, the edge of the pre-specified
grid, which was selected.  On the 30 held-out replicates it gave 0.87 overall (0.84 at the earlier and
0.89 at the later read count), 1.00 for five containers and 0.20 for the near-dead one
(Fig.~\ref{fig:calib}).  That selection has two defects: it used prior width $s{=}1$ while every policy
run uses $s{=}0.5$, and the selected value lies at the grid boundary.

A secondary analysis repeats the selection with $s{=}0.5$ and the extended grid
$\{0,0.003,0.01,0.03,0.1,0.2,0.3,0.5\}$ (Table~\ref{tab:r4calib}).  It re-runs the selection procedure
only.  Under the prior actually used, the best value is $\rho_U{=}0.03$ (selection coverage 0.81, held
out 0.83) rather than 0.1 (0.81 and 0.82), a difference within noise.  Extending the grid does not
rescue calibration: coverage flattens near 0.80 for $\rho_U\in[0.03,0.3]$ and falls to 0.60 at 0.5, so
no value in this likelihood family reaches the 0.90 target.  At every $\rho_U$ the five healthy
containers are covered 0.92--1.00 of the time and the near-dead container 0.00--0.03 of the time.
Widening a likelihood cannot fix a container whose behaviour the model cannot express.  The complete
policy comparison at the re-selected value is reported in Section~\ref{M-sec:res-alloc} of the main paper
(Table~\ref{M-tab:r6lik}).
\begin{figure*}[t]
\centering
\includegraphics[width=0.98\textwidth]{fig_dt_calibration.pdf}
\caption{Is the posterior over container efficiency calibrated?  Each point or bar is the fraction of
cases in which the 90\% credible interval for $\theta_j$ contained the target $\theta^\ast_j$; a
calibrated posterior would sit on the dotted line at 0.90.  \textbf{(a)} Selection replicates
1000--1009, one curve per container, against the beta-binomial correlation $\rho_U$; the thick black
curve is the mean over containers, which is what the selection rule maximised.  \textbf{(b)} Held-out
replicates 0--29 at the selected $\rho_U{=}0.1$, per container.  Five containers are covered perfectly
and the near-dead one is covered in 0.20 of cases at this $\rho_U$ (0.00--0.03 across the extended grid
of Table~\ref{tab:r4calib}).}
\label{fig:calib}
\end{figure*}

%=====================================================================
\section{Completion forecast}\label{supp:forecast}
%=====================================================================
The controller also forecasts the probability of finishing within a remaining budget.  The
\emph{analytic} forecast multiplies per-stripe normal approximations of
$\Prob(\sum_{i\in s}X_i\ge\nu_s)$ with the remaining budget spread evenly over containers.  The
\emph{posterior-simulation} forecast draws 200 values of $\theta$ from the posterior, spreads the
remaining budget in proportion to the reads spent so far, simulates per-fragment recoveries and reports
the fraction of draws in which every stripe completes.  It was evaluated with a budget equal to the
median coverage-greedy cost, at 30/50/70/85\% of it, by expected calibration error (ECE, 10 bins), Brier
score against a constant (climatology) forecast~\cite{gneiting2007strictly} and AUC.

No policy depends on the forecast; we report it because it was pre-specified.  On 1{,}127 checkpoints
from 300 baseline retrievals with a budget of 3.75k reads, of which 52\% finished within the budget, both
forecasts were over-confident (mean forecast 0.86 against 0.52 observed at 30\% of the budget).  The
posterior-simulation forecast beat the analytic one on every metric (ECE 0.29 vs.\ 0.34; Brier 0.30, 95\%
CI 0.27--0.34, vs.\ 0.34, 0.29--0.38; AUC 0.73, 0.69--0.78, vs.\ 0.69, 0.65--0.73), yet both Brier scores
are worse than the 0.25 of a constant forecast at the observed base rate, so neither beats climatology.
They rank retrievals usefully only late in a retrieval (AUC 0.82 at 85\% of the budget).  Futility
stopping, which would rely on them, is disabled in every run, and the forecasts should not be used as
probabilities.
\begin{figure}[t]
\centering
\includegraphics[width=0.6\columnwidth]{fig_forecast.pdf}
\caption{Reliability of the completion forecast on the digital twin: forecast probability of finishing
within the remaining budget against the observed frequency of finishing, pooled over checkpoints at
30/50/70/85\% of the budget (10 equal-width bins, bins with fewer than five checkpoints dropped).  A
calibrated forecast would lie on the diagonal; both curves lie below it, so both are over-confident, and
both have a worse Brier score than a constant forecast at the base rate.}
\label{fig:forecast}
\end{figure}

%=====================================================================
\section{Severity sweep}\label{supp:sweep}
%=====================================================================
The severity sweep (Fig.~\ref{fig:sweep}) varies only the near-dead container's copies per strand.  At
200 and 30 copies, coverage-greedy and Thompson sampling are 2--3\% cheaper than round-robin (2.65k vs.\
2.72k at 200 copies) and tail-aware is even; at 10 copies the policies are within $\pm$1\% except
tail-aware ($+$5\%); at 3 and 1 copies all adaptive rules are costlier (at one copy: coverage-greedy
4.69k, Thompson 4.51k, tail-aware 4.47k vs.\ 4.16k).  Success was 1.00 throughout.  The point at 3 copies
reproduces the baseline design with other seeds (round-robin 3.46k vs.\ 3.49k).  Adaptive allocation
therefore helps only while the model's single multiplier is a fair description of the worst container.
\begin{figure}[t]
\centering
\includegraphics[width=0.6\columnwidth]{fig_dt_sweep.pdf}
\caption{Severity sweep: mean reads consumed as the degraded container's copy number falls from 200 to 1
(300 retrievals per point and policy, 95\% cluster-bootstrap intervals; the axis runs from mild to severe
degradation).  Every retrieval at every level returned the correct file, so the success probability is
1.00 throughout and is not plotted.  Adaptive rules are cheaper than round-robin only at the mild end.}
\label{fig:sweep}
\end{figure}

%=====================================================================
\section{Robustness details}\label{supp:robust}
%=====================================================================
\textbf{Prior assumptions.}  The controller's generic assumptions were varied one at a time on the
baseline (30 syntheses $\times$ 3 read orders each): the prior probability of outright failure
($10^{-2}$, $10^{-1}$ against the primary $10^{-3}$), the prior dropout (0, 5\% against 0.5\%), the
negative-binomial shape (3, 30 against 10), the $\theta$ grid (40 points over $[0.001,10]$ and 10 points
over $[0.01,5]$ against 25 over $[0.01,5]$), the prior width ($s{=}0.25$, $1.0$ against $0.5$) and the
consensus-success curve (scaled by 0.5 and 2.0).  Under the one-parameter model, coverage-greedy stays
between $+$10.9\% and $+$20.6\% against round-robin in every setting except the tighter prior
$s{=}0.25$, where it rises to $+$56.4\%.  Changing the failure prior by two orders of magnitude has no
effect ($+$17.2\% either way), because the failure atom is not what is mismodelled.  Under the
two-parameter model the same sweep stays between $-$0.3\% and $+$3.4\%.

\textbf{Decoder configuration.}  Table~\ref{tab:r4cap} varies the consensus cap.
\input{tables/r4_cap}

\textbf{Failure modes beyond total loss.}  Table~\ref{tab:r4fail} gives the four additional failure
modes summarised in Section~\ref{M-sec:res-robust}.
\input{tables/r4_fail}

\textbf{Other archive geometries.}  Table~\ref{tab:r4geom} gives the six additional geometries.
\input{tables/r4_geoms}

%=====================================================================
\section{Supporting channel result (real nanopore reads)}\label{supp:nanopore}
%=====================================================================
On two independent nanopore flowcells of DNAformer reads~\cite{barlev2025deep,technion2025dnaformerreads},
a strand-specific error propensity persists across flowcells ($r{=}0.57$; null with matched contexts
$0.002$), and a hierarchical beta-binomial error model fitted on one flowcell predicts consensus success
on held-out strands of the other with mean absolute error 0.083, against 0.181 for an i.i.d.\ model.  This
motivates over-dispersed models but does not enter the controller.

%=====================================================================
\section{Upper bound with evaluation leakage}\label{app:leak}
%=====================================================================
Our first analysis of the real read streams (Table~\ref{tab:leak}) used a ``calibrated'' belief whose
per-strand consensus curves were pooled over reads that included the evaluation split.  The controller
therefore had information about the reads it was evaluated on, and the static plans were tuned on them
too.  We keep the table only as an upper bound with evaluation leakage; no claim in the main paper uses
it.  In the analysis of Table~\ref{M-tab:replay} the primary belief is the generic prior, and the optional
calibrated belief uses only the first 300k reads of each run while evaluation uses the next 900k.  On
Twist, the leaky analysis put coverage-greedy at 55.7k reads and the oracle-tuned static plan at 63.1k;
the analysis without leakage gives 56.7k and 65.7k.  The leak therefore flattered the adaptive policy by
about 1k reads and the static baseline by about 2.6k, and the gap without leakage is larger, not smaller.
\input{tables/leak_appendix}

%=====================================================================
\section{Protocol deviations and post-hoc analyses}\label{app:devfull}
%=====================================================================
\textbf{Deviations.}  The protocol and every change made after it was fixed are recorded in
\texttt{dt\_protocol.yaml}.  (i) The planned CustomArray scenario was replaced, before any policy run, by
DT4DDS's photolithographic synthesis model, because the CustomArray settings carry a 117-nt
position-bias table that cannot be applied to 136-nt strands.  (ii) We suspected from commit timestamps
that the first Twist replay had run with $\rho_U{=}0$ and re-ran it; a direct test showed that it had
used $\rho_U{=}0.1$, and the re-run reproduced it exactly.  (iii) One cached homogeneous replicate was
truncated by an interrupted job; it was regenerated from its seed and the scenario re-run in full.
(iv) The photolithographic scenario was first run with 10 syntheses $\times$ 3 read orders for compute
reasons; the first 50 retrievals of the full run were discarded unseen.  Round-robin was later re-run on
it at the full 30 $\times$ 10 design, with the same outcome (0 of 300 recovered).  (v) The static
allocator crashed when only zero-weight (failed) containers remained; it now reads them out, which can
only affect failed retrievals of the failure scenarios.  The $\rho_U$ rule was applied as specified even
though the selected value lies at the edge of the grid.  (vi) The sequential Genscript replay finished
information-greedy after about six hours but then crashed while computing its paired difference, because
some trials never completed (infinite cost); nothing was saved.  We fixed the bug, ran the other Genscript
conditions in parallel processes with the same per-condition seeds (they reproduced every value logged by
the sequential run), and re-started information-greedy; it had not finished when this version was built.
In the sequential run its 200 trials finished and more than 1\% of them never recovered the file within
the available reads, which is why the 0.99 budget was infinite.  We treat the condition as
\emph{censored}: it is excluded from every average and comparison.  (vii) The secondary analyses of
\texttt{dt\_protocol\_rev4.yaml} were specified before they were run: the dropout-aware container model,
the $2\times2\times2$ ablation, the $\rho_U$ re-selection at $s{=}0.5$ on an extended grid, prior
sensitivity, four further failure modes, six further archive geometries, consensus caps of 12 and 48,
finer and calibration-estimated static families, an extended baseline replicate set, the full-design
photolithographic run, the replay calibration-cost experiment and the cost model.  (viii) The following analyses
are exploratory: the complete comparison at $\rho_U{=}0.03$, the reservation
baselines, and the per-read stop, successive-elimination and yield-threshold baselines on regenerated
streams (Section~\ref{M-sec:res-extra}).  (ix) Every cluster-bootstrap interval was
recomputed with one shared, order-independent resampling scheme (Section~\ref{M-sec:eval-stats}); point
estimates are unchanged, and interval limits moved by at most 0.8 percentage points on the baseline,
1.3 in the failure scenarios and 3.6 for information-greedy.

\textbf{Post-hoc analyses} (not pre-specified).  With the binomial likelihood ($\rho_U{=}0$), on the same
300 baseline retrievals, coverage-greedy used 3.34k reads ($-$4.4\% vs.\ round-robin, against $+$15.1\%
at $\rho_U{=}0.1$), Thompson sampling 3.37k ($-$3.5\%, against $+$7.4\%), tail-aware 3.37k ($-$3.5\%,
against $+$12.9\%) and pivotal 3.64k ($+$4.4\%, against $+$49.3\%); q99 was 4.95--5.70k, still not below
round-robin's 5.15k.  With one failed container the binomial likelihood helped tail-aware ($-$25\% vs.\
$+$2\%) and Thompson sampling ($-$20\% vs.\ $-$9\%) and hurt coverage-greedy and pivotal slightly
($-$15\% vs.\ $-$17\%; $-$8\% vs.\ $-$15\%).  Success and safety were unchanged (no wrong output).  On the
Twist replay, the binomial likelihood made coverage-greedy and Thompson sampling costlier (57.6k each,
against 56.7k and 56.9k at $\rho_U{=}0.1$), the opposite of the digital-twin result: the two settings
disagree on which likelihood allocates better.

%=====================================================================
\section{Claim-to-evidence map}\label{supp:evidence}
%=====================================================================
\input{tables/evidence}

\bibliographystyle{IEEEtran}
\bibliography{references}

\end{document}
